\subsubsection{Per-Graph Findings}

\paragraph{The module graph.}
\label{sec:module_conclusion}
\module{Mathlib}'s module graph is not a random network or a simple hierarchy; it is a \textit{deep, engineered backbone} supporting a vast, shallow application layer. Three features stand out:

\begin{enumerate}
    \item \textit{Cognitive redundancy as a feature.} The 17.5\% edge redundancy (\S\ref{sec:transitive-reduction}) represents a \textit{cognitive API}: developers import umbrella modules for navigability, sacrificing logical minimality for usability. Future tools should distinguish \textit{dependency for compilation} (minimality) from \textit{dependency for discovery} (redundancy).

    \item \textit{Funnel topology and compilation bottlenecks.} The layer-width profile reveals a massive leaf layer ($>1{,}400$ modules) resting on a deep, narrow foundation (DAG depth~$153$). Changes in the narrow waist trigger recompilation cascades, identifying a scalability challenge that may eventually require federated compilation tiers.

    \item \textit{Multidimensional centrality.} In-degree identifies \textit{vocabulary}, PageRank identifies the \textit{axiomatic core}, and betweenness identifies \textit{structural bridges} (\S\ref{sec:sa-centrality}). Refactoring a high-betweenness module is riskier than refactoring a high-in-degree module.
\end{enumerate}

\noindent
These observations suggest that \module{Mathlib} has evolved into a \textit{curated small-world network}: new modules cluster within directories (high local clustering) while bridge modules maintain global connectivity. The module tree $T$ and the import graph $G_{\mathrm{module}}$ continuously reshape each other---developers import along the cognitive contours of $T$, while cross-directory dependencies force periodic reorganization of the naming hierarchy. A detailed discussion is provided in Appendix~\ref{app:module-analysis}.

\paragraph{The declaration graph.}
\label{sec:thm-conclusion}
The declaration graph dissolves the module boundary, exposing individual logical dependencies. Three findings stand out:

\begin{enumerate}
    \item \textit{Double infrastructure.} The hub structure decomposes into a \emph{language infrastructure} layer (\decl{OfNat.ofNat}, \decl{Eq.refl}, \decl{DFunLike.coe}) determined by Lean's type system, and a \emph{mathematical infrastructure} layer (\decl{CategoryTheory.Category}, \decl{Real}, \decl{TopologicalSpace}) determined by mathematical logic (Remark~\ref{rem:two-layers}). The language layer is a process-trace---a different proof assistant would produce a different hub ranking.

    \item \textit{Emergent disciplinary structure.} Louvain community detection recovers seven major disciplines from dependency topology alone (modularity $0.48$; \S\ref{sec:sa-community}), with only moderate alignment to namespaces (NMI $= 0.34$). Only \ns{CategoryTheory} achieves near-complete community--namespace correspondence ($97.4\%$).

    \item \textit{Deep cross-disciplinary entanglement.} Cross-namespace edge density ($50.9\%$) far exceeds the module-level figure ($37.1\%$). Module boundaries conceal the true extent of mathematical interdependence.
\end{enumerate}

\noindent
Even at this finer resolution, the process leaves traces: the language infrastructure layer records Lean's design choices, the $44.8\%$ leaf theorem rate includes \texttt{@[to\_additive]} mirrors, and bounded out-degree encodes the cognitive ceiling of human proof complexity. A detailed discussion and a module-vs-declaration comparison table are provided in Appendix~\ref{app:decl-analysis}.

\paragraph{The namespace graph.}
\label{sec:ns-analysis}
The namespace graph is the only level admitting a continuous parameterization (Definition~\ref{def:ns-graph}), yielding a family $G_{\mathrm{ns}}^{(k)}$ indexed by truncation depth. Three findings stand out:

\begin{enumerate}
    \item \textit{Rapid containment decay.} The organizing power of human-imposed boundaries diminishes steeply with granularity: from $48.9\%$ at the discipline level to ${\sim}14\%$ by depth~$2$, with a ``knee'' at the discipline--sub-discipline transition costing $26.7$ percentage points in a single step.

    \item \textit{Structural interpolation.} $G_{\mathrm{ns}}^{(2)}$ consistently interpolates between the module and declaration levels: modularity ($0.283$) and NMI ($0.253$) are both lower than at the declaration level, reflecting the coarsening effect of namespace aggregation.

    \item \textit{Infrastructure bottleneck.} The namespace graph is more vulnerable to targeted attack than the declaration graph: $20\%$ targeted removal reduces the GCC to $20.8\%$, versus $46.1\%$ at the declaration level (\S\ref{sec:sa-robustness}). This heightened fragility reveals the structural dependence of formalized mathematics on a narrow foundation of shallow infrastructure.
\end{enumerate}

\noindent
The containment decay curve is a quantitative measure of a cultural artifact's explanatory power: the disciplinary taxonomy ``works'' at the broadest scale but breaks down at sub-discipline granularity. A detailed discussion is provided in Appendix~\ref{app:ns-analysis}.
